\documentclass[11pt,a4paper]{article}
\usepackage[main=italian,provide=*]{babel}
\usepackage[utf8]{inputenc}
\usepackage[T1]{fontenc}
\usepackage{amsmath,amssymb}
\usepackage{geometry}
\geometry{margin=2.7cm}
\usepackage{booktabs}
\usepackage{array}
\usepackage{caption}
\usepackage[colorlinks=true,linkcolor=blue,citecolor=blue,urlcolor=blue]{hyperref}

\title{Quantum simulation del dimero: applicazione della teoria\\
all'implementazione Qiskit\\[6pt]
\large Dai notebook di teoria ai notebook di simulazione --- procedura commentata}
\author{Samuele \\ Relatore: Prof. Alessandro Chiesa}
\date{20 luglio 2026}

\begin{document}
\maketitle

\begin{abstract}
\noindent Questo documento ricostruisce, passo per passo, come la teoria
raccolta in \texttt{quantum\_simulation\_dimero\_teoria.md} e nei documenti
di derivazione del dimero sia stata applicata nella costruzione dei due
notebook Qiskit prodotti per il task di \emph{quantum simulation} (evoluzione
temporale $e^{-iHt}$) richiesto dal relatore:
\texttt{quantum\_simulation\_dimero\_trotter.ipynb} (versione completa, con
derivazione fisica del punto di lavoro) e la sua versione
\texttt{\_semplificato.ipynb} (solo il richiesto dal relatore, con
esplorazione interattiva dei parametri). Ogni passaggio riporta il
riferimento alla teoria da cui discende; dove la teoria di progetto non
copre un passaggio, si cita la fonte esterna usata. Le versioni dei
documenti di teoria citate sono quelle \textbf{corrette} (vedi
\texttt{errata\_simmetrie\_dmi\_spin} e la ricostruzione integrale di
\texttt{simmetrie\_dmi\_spin.tex}, oltre alla correzione minore di
\texttt{dimero\_sintesi.tex}).
\end{abstract}

\tableofcontents

\section{Mappa dei riferimenti}

La tabella seguente elenca i documenti di teoria richiamati in questo
documento e il loro stato.

\begin{center}
\begin{tabular}{@{}p{4.6cm}p{9.6cm}@{}}
\toprule
Documento & Stato e uso \\
\midrule
\texttt{quantum\_simulation\_dimero\_teoria.md} & Derivazione completa dei blocchi
esatti di Trotter (\S2--\S6 di quel file). Nessuna correzione necessaria: è
il riferimento primario di questo documento. \\
\texttt{quantum\_simulation\_trascrizione\_appunti.md} & Trascrizione degli
appunti del relatore sull'Hamiltoniana $H=H_1+H_2$; riferimento per
l'inquadramento iniziale (\S1). \\
\texttt{dimero\_DM.pdf} (alias \texttt{dimero\_teoria.pdf}) & Derivazione
completa dello spettro a matrice $4\times4$, blocco $3\times3$, autostato
esatto $|t_0\rangle$. Verificato interamente corretto in questa
conversazione: nessuna correzione. Usato per \S2 e \S5. \\
\texttt{simmetrie\_dmi\_spin.tex} & Regole di selezione della DMI, struttura
a V, isotropia del singoletto. \textbf{Versione corretta} (fattore $2$ e
$\sqrt2$ negli elementi di matrice, propagato a 8 sezioni): usata per \S2. \\
\texttt{dimero\_sintesi.tex} & Anti-incrocio a $b/J=2$, degenerazione a
$B=0$. \textbf{Corretta} in un solo punto (norma dei commutatori in \S2.2 di
quel file): non direttamente citata qui ma concettualmente affine a \S5. \\
\texttt{dimero\_spin.pdf} & Origine microscopica $J=4t^2/U$ da Hubbard.
Verificato corretto; non necessario in questo documento (riguarda l'origine
di $J$, non la simulazione Trotter). \\
\bottomrule
\end{tabular}
\end{center}

Dove non esiste un riferimento nei documenti di progetto, il documento cita
esplicitamente la fonte esterna (letteratura o documentazione Qiskit) usata
per colmare il vuoto.

\section{Punto di partenza: la richiesta del relatore}

\subsection{La richiesta}

Il relatore ha chiesto (risposte del 19/07/2026, vedi
\texttt{domande\_relatore.md} e \texttt{log\_decisioni.md}) di esplorare
diversi set di parametri con la dinamica \emph{esatta} per trovare
oscillazioni di $\langle S_z\rangle$ non banali (``non monocromatiche''),
partendo dal set $J=2b$, $D=b/3$. Il compito è quindi in due fasi distinte:
(i) trovare con la dinamica esatta un punto di lavoro che soddisfi il
requisito qualitativo; (ii) implementare l'evoluzione via Suzuki--Trotter su
quel punto. Questo documento segue esattamente quest'ordine.

\subsection{Perché serve il termine DM: riferimento alla teoria}

L'Hamiltoniana di riferimento, confermata dal relatore e trascritta in
\texttt{quantum\_simulation\_trascrizione\_appunti.md},
$$H = b(s_{z1}+s_{z2}) + J\,\vec s_1\cdot\vec s_2 + D(s_{x1}s_{z2}-s_{z1}s_{x2}) = H_1+H_2,$$
separa un blocco $H_1$ (campo + scambio, fisso) da un blocco $H_2$ (solo DM,
si annulla per $D=0$). Il documento \texttt{dimero\_DM.pdf} (Tabella 1, limite
$D\to0$) mostra che lo stato $|t_+\rangle=|00\rangle$ è autostato esatto di
$H_1$ con $E=J+2b$, per qualunque $b,J$. Poiché $S_z|00\rangle=|00\rangle$ è
anch'esso autostato (a autovalore $+1$), segue che $\langle S_z\rangle(t)$ è
identicamente costante quando $D=0$: \textbf{è il termine DM, e solo esso, a
rendere possibile un'oscillazione}. Questo fissa il vincolo operativo: il
punto di lavoro cercato deve avere $D\neq0$.

\subsection{Il set suggerito fallisce: verifica numerica}

Applicando la dinamica esatta $e^{-iHt}|00\rangle$ al set suggerito
($b/J=0.5$, $D/J=1/6$, chiamato $R_0$ nei notebook), l'escursione
picco-picco di $\langle S_z\rangle(t)$ risulta $\approx0.0062$: sotto la
soglia del rumore statistico di un campionamento a $8192$ shot
($1/\sqrt{8192}\approx0.011$). Il set $R_0$ non soddisfa il requisito, e va
sostituito con un punto derivato dalla struttura fisica del problema --
oggetto della sezione seguente.

\section{Derivazione fisica del punto di lavoro $R_1$}

\subsection{Passaggio alla base di spin totale}

Il documento \texttt{dimero\_DM.pdf} (Passo 3, eq.~6--7) introduce il cambio
di base
$$|t_0\rangle=\frac{|01\rangle+|10\rangle}{\sqrt2}, \qquad |s\rangle=\frac{|01\rangle-|10\rangle}{\sqrt2}, \qquad |t_\pm\rangle=|00\rangle,|11\rangle,$$
in cui $H_1$ (campo + scambio) è diagonale:
$$E(t_+)=J+2b,\quad E(t_0)=J,\quad E(s)=-3J,\quad E(t_-)=J-2b.$$
Questa è esattamente la base $\{|T_{+1}\rangle,|T_0\rangle,|S\rangle,|T_{-1}\rangle\}$
di \texttt{simmetrie\_dmi\_spin.tex} (versione corretta), a meno della
convenzione di normalizzazione dell'energia ($J$ del notebook Trotter
corrisponde a $4J$ della convenzione $s=\sigma/2$ di quel documento --
coerente con la distinzione $J_{\rm code}$ vs $J_{\rm Crippa}$ già tracciata
nel resto del progetto).

\subsection{Struttura a V degli accoppiamenti DM}

Il documento \texttt{simmetrie\_dmi\_spin.tex} (versione corretta, \S11.1 e
\S11.2.2) stabilisce, tramite l'argomento di simmetria con l'operatore di
scambio $P_{12}$, che il termine DM:
\begin{itemize}
\item non accoppia mai il singoletto con se stesso, né i tripletti fra loro;
\item accoppia il singoletto a \emph{entrambi} i tripletti $M=\pm1$ con la
stessa forza (Caso B, $\vec D$ trasversale -- la nostra situazione, poiché
$H_2\propto s_{x1}s_{z2}-s_{z1}s_{x2}\propto(\vec s_1\times\vec s_2)_y$);
\item lascia $|T_0\rangle$ completamente disaccoppiato.
\end{itemize}
Con il fattore corretto (Errata di \texttt{simmetrie\_dmi\_spin.tex},
\S11.2.2), l'accoppiamento vale $V=D_\perp/(2\sqrt2)$ nella convenzione
$s=\sigma/2$. Riportato alla convenzione del notebook Trotter (dove il
termine è scritto come $D(s_{x1}s_{z2}-s_{z1}s_{x2})$ con lo stesso
significato di $D$), la struttura è verificata numericamente nel notebook
stesso: il DM accoppia $|00\rangle=|t_+\rangle$ e $|11\rangle=|t_-\rangle$ a
$|s\rangle$ con la stessa forza $\sqrt2 D$ (coerente con l'eq.~14 di
\texttt{dimero\_DM.pdf}, $g=\sqrt2D$), e $|t_0\rangle$ risulta esattamente
disaccoppiato -- verificato sia analiticamente (\texttt{dimero\_DM.pdf},
eq.~10) sia con il calcolo diretto nel notebook.

\`E dunque un \textbf{sistema a V}: $|t_+\rangle\leftrightarrow|s\rangle\leftrightarrow|t_-\rangle$,
con lo stato iniziale $|00\rangle=|t_+\rangle$ su un braccio. I due parametri
di controllo sono i detuning
$$\Delta_+=E(t_+)-E(s)=J+2b, \qquad \Delta_-=E(t_-)-E(s)=J-2b$$
(nella normalizzazione del notebook Trotter, dove l'accoppiamento
effettivo è $V\propto D$).

\subsection{Prima condizione: ampiezza del segnale (formula di Rabi)}

Nessuno dei documenti di teoria di progetto deriva la formula di
popolazione trasferita per un sistema a due livelli con detuning e
accoppiamento generici: è un risultato standard di meccanica quantistica,
non specifico al progetto. Riferimento esterno usato:
\href{https://web.pa.msu.edu/people/mmoore/Lect7_RabiModel.pdf}{M.\ Moore,
\emph{General Study of Two-Level Systems}, dispense del corso, Michigan
State University} -- vi si deriva la popolazione massima trasferita in un
sistema a due livelli con accoppiamento $\Omega$ (qui $V$) e detuning
$\Delta$:
$$P_{\max} = \frac{4V^2}{4V^2+\Delta^2}.$$
Applicando questa formula al braccio $|t_+\rangle\leftrightarrow|s\rangle$
del sistema a V (accoppiamento $V\propto D$, detuning $\Delta_+=J+2b$ nella
normalizzazione del notebook, dove risulta $V=D$ dopo la riduzione al
sottoproblema 2$\times$2 -- coerente con l'eq.~19 di
\texttt{dimero\_DM.pdf} per il caso della risonanza $b/J=2$, che è lo stesso
tipo di sottoproblema 2$\times$2), per $V\ll\Delta_+$:
$$\text{escursione di }\langle S_z\rangle \;\simeq\; \frac{4V^2}{\Delta_+^2} \;=\; \frac{D^2}{2(b+J)^2}$$
(fattore di normalizzazione verificato empiricamente nel notebook: la
formula riproduce l'escursione esatta di $R_0$ a quattro cifre, $0.0062$
contro $0.0062$). \textbf{Conclusione procedurale:} $R_0$ fallisce perché sta
in pieno regime perturbativo ($D/J=1/6\ll1$); serve $D$ dell'ordine di $J$.

\subsection{Seconda condizione: numero di frequenze}

Nessun documento di teoria di progetto tratta esplicitamente il criterio
``non monocromatico'' in termini di frequenze di Bohr multiple; è
un'elaborazione originale di questo lavoro, basata sulla decomposizione
spettrale standard $\langle S_z\rangle(t)=\sum_k|c_k|^2(S_z)_{kk}+\sum_{k<l}2\mathrm{Re}[c_k^*c_l(S_z)_{kl}e^{-i\Delta_{kl}t}]$
(teoria elementare di meccanica quantistica, non richiede citazione
specifica). Il campo $b$ controlla quante frequenze di Bohr compaiono con
ampiezza confrontabile:
\begin{itemize}
\item $b\to-J$ (risonanza, $\Delta_+\to0$): il braccio $t_+$--$s$ domina,
riducendo il sistema a un due-livelli -- oscillazione di Rabi ad ampiezza
piena ma \textbf{monocromatica} (verificato: $a_2/a_1\to0$);
\item $b\to0$ (bracci simmetrici, $\Delta_+=\Delta_-$): un solo modo
("bright") domina di nuovo.
\end{itemize}
L'ottimo sta a $|b|$ piccolo ma non nullo. Uno scan 1D (a $D/J=1$ fissato
dalla prima condizione) localizza il massimo del rapporto fra le ampiezze
dei due modi dominanti a $b/J\simeq-0.18$.

\subsection{Punto adottato}

$$\boxed{R_1 = (b/J=-0.18,\; D/J=1.0)}$$
con escursione $0.84$ e $a_2/a_1=0.995$ (due frequenze dominanti di
ampiezza praticamente identica). Verificata la robustezza: tutta la fascia
$b/J\in[-0.22,-0.14]$ soddisfa entrambe le condizioni, quindi la posizione
esatta del massimo non ha significato fisico speciale -- conta la regione.
Il segno negativo di $b/J$ non è artificioso: sotto flip globale
$\Pi=X\otimes X$ vale $\Pi H(b,J,D)\Pi=H(-b,J,-D)$, ma $\Pi|00\rangle=|11\rangle$,
e lo stato iniziale $|00\rangle$ è fissato dal relatore.

\subsection{Nota procedurale: due versioni scartate}

Il punto $R_1$ è stato raggiunto dopo due tentativi precedenti, entrambi
scartati e documentati per trasparenza procedurale (vedi
\texttt{quantum\_simulation\_dimero\_trotter\_note.md}):
\begin{enumerate}
\item una prima versione usava uno scan automatico su 2349 punti classificati
con un'\emph{entropia spettrale pesata} non riconducibile ad alcuna fonte di
progetto -- rimossa per assenza di giustificazione bibliografica e per bassa
robustezza (scarto 1\textsuperscript{o}--5\textsuperscript{o} classificato
del solo 4\%);
\item una seconda versione usava una tabella di regimi candidati costruita
\emph{a posteriori} con $R_1$ già dentro -- una razionalizzazione, non una
derivazione, che non permetteva di ricostruire come si fosse arrivati al
punto.
\end{enumerate}
La derivazione corrente (\S2.1--2.5) è quella che ha sostituito entrambi i
tentativi, con ogni passaggio riconducibile o alla teoria di progetto
(struttura a V) o a un risultato standard esplicitamente citato (formula di
Rabi).

\section{Decomposizione in gate: applicazione diretta della teoria}

Questa sezione applica letteralmente le formule di
\texttt{quantum\_simulation\_dimero\_teoria.md}, senza modifiche: è la parte
del lavoro in cui la teoria è stata usata come sorgente diretta del codice,
non solo come sfondo concettuale.

\subsection{Blocco di scambio $U_J(t)$}

Da \texttt{quantum\_simulation\_dimero\_teoria.md} \S2: con
$s_i^\alpha=\sigma_i^\alpha/2$, $J\vec s_1\cdot\vec s_2=\frac{J}{4}(XX+YY+ZZ)$,
e poiché $XX,YY,ZZ$ commutano a due a due (\S2.2 della teoria, verificato
con le regole di Pauli $XY=iZ$ etc.), l'esponenziale fattorizza
\textbf{esattamente}:
$$U_J(t) = R_{XX}\!\left(\frac{Jt}{2}\right)R_{YY}\!\left(\frac{Jt}{2}\right)R_{ZZ}\!\left(\frac{Jt}{2}\right).$$
Il fattore $Jt/2$ (anziché $Jt/4$) discende dalla convenzione Qiskit
$R_{XX}(\theta)=e^{-i\frac\theta2 XX}$, verificata numericamente nella
teoria \S2.3 con \texttt{RXXGate(0.7).to\_matrix()}. Nel notebook applicativo
questa stessa verifica è stata ripetuta per il caso specifico $R_1$
($b,J,D$ numerici) come self-test, non come nuova derivazione.

\subsection{Blocco campo+scambio $e^{-iH_1t}$}

Da \texttt{quantum\_simulation\_dimero\_teoria.md} \S3: il campo commuta con
lo scambio isotropo perché $[S^2,S_z]=0$ (Casimir dell'algebra del momento
angolare, un fatto generale di meccanica quantistica non specifico al
progetto ma richiamato esplicitamente nella teoria di riferimento), quindi
$$e^{-iH_1t} = R_z^{(1)}(bt)\,R_z^{(2)}(bt)\,U_J(t),$$
\textbf{esatto}, nessun errore di approssimazione. Questo blocco è
implementato nel notebook come funzione \texttt{step\_H1}.

\subsection{Blocco DM $U_2(t)$}

Da \texttt{quantum\_simulation\_dimero\_teoria.md} \S4: con $P_1=X_1Z_2$,
$P_2=Z_1X_2$, si verifica $[P_1,P_2]=0$ (§4.2 della teoria, tramite
$P_1P_2=P_2P_1=Y_1Y_2$), quindi la fattorizzazione è esatta. Usando
$HZH=X$ (Hadamard scambia $X\leftrightarrow Z$):
$$U_2(t) = \Big[(I\otimes H)R_{ZZ}\!\big(-\tfrac{Dt}{2}\big)(I\otimes H)\Big]\cdot\Big[(H\otimes I)R_{ZZ}\!\big(\tfrac{Dt}{2}\big)(H\otimes I)\Big].$$
Implementato nel notebook come funzione \texttt{step\_H2}, con il commento
esplicito che documenta la nota sul segno (\S4.5 della teoria): se il
termine DM avesse convenzione di segno opposta, entrambi gli angoli
cambierebbero segno insieme, senza alterare la struttura del circuito.

\subsection{Perché i blocchi isolati sono esatti (BCH)}

Da \texttt{quantum\_simulation\_dimero\_teoria.md} \S5: per la formula di
Baker--Campbell--Hausdorff, $[A,B]=0$ annulla \emph{ogni} termine della
torre di commutatori annidati, non solo il primo -- quindi tutte e tre le
fattorizzazioni sopra sono esatte, non approssimazioni al prim'ordine
troncate. Solo la separazione fra $H_1$ e $H_2$ (dove
$[H_1,H_2]\neq0$ per $D\neq0$) introduce l'approssimazione, trattata in
\S4 di questo documento.

\subsection{Self-test numerico dei blocchi}

Applicazione pratica della verifica di convenzione già fatta nella teoria
(\S2.3): nel notebook, ciascun blocco (\texttt{step\_H1}, \texttt{step\_H2},
il circuito completo) è confrontato con l'esponenziale di matrice
(\texttt{scipy.linalg.expm}) a un tempo arbitrario $\tau=0.37$, con errore
atteso a precisione macchina ($\sim10^{-16}$). Verificato: errore blocco
$H_1=2.3\times10^{-16}$, blocco $H_2=2.2\times10^{-16}$, circuito completo
vs matrice ($N=5$) $=8.3\times10^{-16}$. Un errore fuori da questo ordine di
grandezza segnalerebbe un problema di convenzione (fattori $1/2$, segni,
ordinamento dei qubit), non un errore di Trotter -- da qui la necessità di
eseguire questo test \emph{prima} di ogni risultato fisico, come esplicitato
nel notebook.

\section{Errore di Trotter: dalla formula generale alla verifica numerica}

\subsection{Formula dell'errore}

Da \texttt{quantum\_simulation\_dimero\_teoria.md} \S6: per un passo di
durata $\tau=t/N$,
$$U_{\rm esatto}(\tau)-U_{\rm Trotter}(\tau) = \frac{\tau^2}{2}[H_1,H_2]+O(\tau^3),$$
da cui, per la disuguaglianza triangolare su $N$ passi, l'errore totale
scala come $O(t^2/N)$ (\S6.3 della teoria). Conseguenza pratica (\S6.4):
raddoppiando $t$ a $N$ fisso l'errore quadruplica; per tenerlo costante
serve $N\propto t^2$.

\subsection{Verifica di convergenza e scaling}

Applicando questa formula al punto $R_1$: a $t=10$ fissato, la pendenza
log-log dell'infedeltà $1-\mathcal F$ misurata è $-2.02$ (atteso $-2$,
quadratico nell'errore di ampiezza) e quella della norma $\|\Delta U\|_2$ è
$-1.06$ (atteso $-1$, primo ordine) -- entrambe coerenti con la teoria di
\S6.5. A $N$ piccolo l'errore non è monotono ($N=2$ dà $1.17$, $N=5$ dà
$1.44$): fuori dal regime asintotico $\tau\|H\|\ll1$ lo scaling non si
applica, coerentemente con la natura asintotica della formula BCH troncata.

\subsection{Un fenomeno oltre la teoria generale: scaling $N^{-2}$ anomalo sull'osservabile}

Questo è un risultato che \textbf{non discende direttamente} da
\texttt{quantum\_simulation\_dimero\_teoria.md}, e per cui non esiste un
riferimento diretto nei documenti di progetto: l'errore sull'osservabile
$\langle S_z\rangle$ (non sulla fedeltà, ma sulla singola quantità
misurata) scala empiricamente come $N^{-2}$ anziché $N^{-1}$, pur essendo
lo stato-errore lineare in $1/N$ (verificato: $\|\Delta U\|_2\sim N^{-1.06}$).

La spiegazione, sviluppata in questo lavoro: l'Hamiltoniana del dimero in
questa forma è \textbf{reale simmetrica} (il termine DM $X_1Z_2-Z_1X_2$ è
reale), quindi $[H_1,H_2]$ è \textbf{reale antisimmetrica}. Il generatore
d'errore al prim'ordine $\varepsilon=\frac\tau2 i[H_1,H_2]$ è allora
hermitiano ma con \textbf{diagonale identicamente nulla in ogni base reale
ortogonale} (verificato: diagonale $=0$ a precisione macchina nella base di
$H$), quindi le frequenze di Bohr non hanno errore $O(\tau)$: lo spostamento
parte da $O(\tau^2)$ (verificato: dimezzando $\tau$ lo scarto si divide per
$4.00$).

Questo è un caso particolare -- non trattato esplicitamente nella teoria
generale di \S6 -- di un fenomeno più ampio studiato in letteratura: la
struttura di commutatori specifica di un problema (non solo la loro norma)
può ridurre lo scaling dell'errore di Trotter rispetto al caso generico.
Riferimento esterno per approfondimento:
\href{https://arxiv.org/abs/1912.08854}{A.\,M.\ Childs, Y.\ Su, M.\,C.\ Tran,
N.\ Wiebe, S.\ Zhu, \emph{Theory of Trotter Error with Commutator Scaling},
Phys.\ Rev.\ X \textbf{11}, 011020 (2021)}, che sviluppa limiti d'errore più
stretti sfruttando esplicitamente la struttura dei commutatori piuttosto che
il solo troncamento BCH. \textbf{Attenzione:} questo vantaggio è specifico di
questa Hamiltoniana e di questa partizione $H_1/H_2$ (entrambe reali); con
un termine DM in forma complessa (es.\ $XY-YX$) $H$ non sarebbe reale e ci
si aspetterebbe il comportamento $N^{-1}$ generico. Non va generalizzato
implicitamente all'estensione a $N=3$ o alla Parte 2 (rumore) senza rifare
la verifica.

\section{Struttura a 3 livelli: il disaccoppiamento di $|T_0\rangle$}

\subsection{Origine teorica}

Che $|t_0\rangle$ sia esattamente disaccoppiato non è un'osservazione
originale di questo lavoro: discende direttamente da
\texttt{dimero\_DM.pdf} (Passo 4, eq.~10: $H_{\rm DM}|t_0\rangle=0$, e Passo
5, eq.~13: $H|t_0\rangle=J|t_0\rangle$ esatta per ogni $b,D$) ed è coerente
con le regole di selezione di \texttt{simmetrie\_dmi\_spin.tex} (versione
corretta, \S11.1): il termine DM connette solo settori di simmetria di
scambio opposta, e $|t_0\rangle$ (combinazione simmetrica) non ha alcun
canale verso se stesso né verso il singoletto lungo l'asse specifico
coinvolto in questa geometria.

\textbf{Applicazione:} nel notebook, questo fatto teorico è stato verificato
numericamente nel punto $R_1$ specifico (non ri-derivato): $\langle
T_0|00\rangle=0$ (non popolato dallo stato iniziale), $S_z|T_0\rangle=0$
(non contribuisce a $S_z$), $H|T_0\rangle=E\,|T_0\rangle$ con
$E=J/4$ nella convenzione $s=\sigma/2$ -- verificato per tre set di
parametri diversi, confermando l'indipendenza da $b,D$ prevista dalla
teoria.

\subsection{Trotter eredita il disaccoppiamento}

Estensione originale (non presente nella teoria di progetto): poiché
$|T_0\rangle$ è autostato di $H_1$ \emph{e} annichilato da $H_2$
\emph{separatamente}, entrambi i fattori $e^{-iH_1\tau}$ e $e^{-iH_2\tau}$
preservano individualmente il disaccoppiamento -- e quindi anche il loro
prodotto (Trotter). Verificato numericamente: popolazione spuria di
$|T_0\rangle$ sotto evoluzione di Trotter $=0.0000$ per $N=5,20,50,100$,
identica alla dinamica esatta ($1.3\times10^{-30}$). Conseguenza generale:
l'errore di Trotter non è rumore isotropo nello spazio di Hilbert, ma vive
solo nel sottospazio dove i pezzi $H_1,H_2$ non commutano -- coerente con
la Sezione 6 di \texttt{quantum\_simulation\_dimero\_teoria.md}, che lega
l'errore proprio al commutatore $[H_1,H_2]$.

\section{Esecuzione su circuito e campionamento}

Questa parte non ha un riferimento diretto nei documenti di teoria di
progetto (che si fermano alla decomposizione in gate): è ingegneria di
misura standard, per cui si cita la fonte esterna.

$S_z=(Z_1+Z_2)/2$ è diagonale nella base computazionale, con autovalori
$+1$ su $|00\rangle$, $-1$ su $|11\rangle$, $0$ sugli altri due stati.
La stima da conteggi,
$$\langle S_z\rangle \simeq \frac{n_{00}-n_{11}}{N_{\rm shots}},$$
segue direttamente dalla regola di Born per la misura in base computazionale
(si veda ad es.\ la documentazione Qiskit sulle primitive di misura,
\url{https://docs.quantum.ibm.com/guides/get-started-with-primitives}, o
qualunque testo standard di calcolo quantistico per la derivazione generale
della regola di Born). L'errore statistico atteso $\sim1/\sqrt{N_{\rm
shots}}$ è la deviazione standard di una proporzione binomiale.

\textbf{Applicazione e risultato:} nel punto $R_1$ con $N=20$ e $8192$
shot, l'errore di Trotter (statevector vs esatto, $\approx0.25$) risulta
oltre 30 volte il rumore statistico ($\approx0.007$--$0.011$): l'errore
dominante in questo regime è \emph{algoritmico}, non statistico. Aumentare
gli shot non migliora il risultato; serve aumentare $N$ -- una conclusione
procedurale rilevante per la Parte 2 (rumore), dove il compromesso fra
errore di Trotter (che vuole $N$ grande) ed errore di gate fisico (che
vuole $N$ piccolo) diventerà centrale.

\section{Dal notebook completo al notebook semplificato: scelte procedurali}

Questa sezione documenta le decisioni prese nella seconda fase del lavoro,
su richiesta esplicita di ridurre il notebook a ciò che il relatore ha
richiesto, mantenendo la possibilità di esplorare interattivamente i
parametri.

\subsection{Cosa è stato rimosso}

La derivazione fisica del sistema a V (\S2 di questo documento) è stata
tenuta fuori dal notebook semplificato: il notebook completo la mantiene
come giustificazione rigorosa, quello semplificato la sostituisce con
un'esplorazione interattiva (widget) che permette di \emph{osservare}
direttamente le due condizioni (escursione, $a_2/a_1$) senza doverle
derivare analiticamente -- coerente con la richiesta di presentare il
lavoro come proprio, con la ricerca dei parametri resa esplicita e
manipolabile piuttosto che dimostrata a priori.

\subsection{Widget interattivi e persistenza dei parametri}

Implementazione: due \texttt{FloatSlider} (\texttt{ipywidgets}) per $b/J$ e
$D/J$, che ridisegnano in tempo reale la dinamica esatta e lo spettro dei
modi; un pulsante ``Fissa questi parametri'' scrive la scelta in
\texttt{parametri\_scelti.json}, letto dalle sezioni successive del
notebook. Nessun riferimento di teoria necessario: è una scelta di
ingegneria del notebook per rendere esplicito, riproducibile e persistente
il passaggio dalla ricerca esplorativa al valore adottato. Fallback esplicito
se \texttt{ipywidgets} non è disponibile (chiamata diretta della funzione
\texttt{esplora(b/J, D/J)}).

\subsection{Correzioni tecniche}

Due difetti individuati e corretti nel corso dello sviluppo, non
riconducibili a teoria fisica ma a corretta gestione delle risorse
matplotlib/Qiskit:
\begin{itemize}
\item \textbf{Duplicazione delle immagini}: ogni chiamata alle funzioni di
plotting creava una nuova \texttt{Figure} senza chiuderla
(\texttt{plt.close()}), causando accumulo di figure aperte che potevano
essere ridisegnate dal flush automatico di Jupyter. Verificato: 5 chiamate
senza chiusura lasciavano 5 figure aperte; con \texttt{plt.close(fig)} dopo
ogni \texttt{plt.show()}, 0 figure aperte anche dopo chiamate ripetute.
\item \textbf{Punti campionati non coerenti con la selezione}: nella figura
riassuntiva a curve selezionabili, i punti da campionamento (shots) erano
calcolati una sola volta per $N=20$ e riusati per qualunque selezione di
$N$ nel menu. Corretto con una cache (\texttt{\_cache\_shots}) che calcola i
punti per ciascun $N$ effettivamente selezionato, la prima volta che viene
richiesto, e li riusa alle chiamate successive.
\end{itemize}

\section{Riepilogo dei riferimenti per sezione}

\begin{center}
\begin{tabular}{@{}p{2.4cm}p{5.6cm}p{6.2cm}@{}}
\toprule
Sezione di questo doc.\ & Riferimento di teoria (progetto) & Riferimento esterno (dove serve) \\
\midrule
\S1 & \texttt{trascrizione\_appunti.md}; \texttt{dimero\_DM.pdf} Tab.\,1 & --- \\
\S2.1--2.2 & \texttt{dimero\_DM.pdf} Passo 3; \texttt{simmetrie\_dmi\_spin.tex} corretto \S11.1, \S11.2.2 & --- \\
\S2.3 & --- & MSU, \emph{Rabi Model} (formula $4V^2/(4V^2+\Delta^2)$) \\
\S2.4--2.6 & --- (elaborazione originale) & --- \\
\S3 & \texttt{quantum\_simulation\_dimero\_teoria.md} \S2--\S5 (per intero) & Qiskit \texttt{RXXGate} docs (verifica \S2.3 teoria) \\
\S4.1--4.2 & \texttt{quantum\_simulation\_dimero\_teoria.md} \S6 & --- \\
\S4.3 & --- (estensione originale) & Childs et al., PRX 11, 011020 (2021) \\
\S5 & \texttt{dimero\_DM.pdf} Passo 4--5; \texttt{simmetrie\_dmi\_spin.tex} corretto \S11.1 & --- \\
\S6 & --- & Qiskit measurement primitives docs; regola di Born (standard) \\
\S7 & \texttt{quantum\_simulation\_dimero\_trotter\_note.md} (cronologia $R_1$) & --- \\
\bottomrule
\end{tabular}
\end{center}

\section*{Riferimenti bibliografici completi}

\begin{itemize}
\item H.\,F.\ Trotter, \emph{On the product of semi-groups of operators},
Proc.\ Am.\ Math.\ Soc.\ \textbf{10}, 545--551 (1959).
\item M.\ Suzuki, \emph{Generalized Trotter's formula and systematic
approximants of exponential operators and inner derivations with
applications to many-body problems}, Commun.\ Math.\ Phys.\ \textbf{51},
183--190 (1976).
\item S.\ Lloyd, \emph{Universal quantum simulators}, Science \textbf{273},
1073--1078 (1996).
\item A.\,M.\ Childs, Y.\ Su, M.\,C.\ Tran, N.\ Wiebe, S.\ Zhu, \emph{Theory
of Trotter Error with Commutator Scaling}, Phys.\ Rev.\ X \textbf{11},
011020 (2021); arXiv:1912.08854.
\item M.\ Moore, \emph{General Study of Two-Level Systems} (dispense del
corso, Michigan State University),
\url{https://web.pa.msu.edu/people/mmoore/Lect7_RabiModel.pdf}.
\item IBM Quantum, \emph{Get started with primitives} (documentazione
Qiskit), \url{https://docs.quantum.ibm.com/guides/get-started-with-primitives}.
\item M.\ Crippa et al., \emph{Simulating Static and Dynamic Properties of
Magnetic Molecules with Prototype Quantum Computers}, Magnetochemistry
\textbf{7}, 117 (2021) -- riferimento concettuale generale del progetto.
\end{itemize}

\appendix
\section{File prodotti in questa fase del lavoro}

\begin{itemize}
\item \texttt{quantum\_simulation\_dimero\_trotter.ipynb} -- notebook
completo, con derivazione fisica di $R_1$, riferimenti interni alla
teoria, sezione errata/discrepanza con \texttt{simmetrie\_dmi\_spin.pdf}.
\item \texttt{quantum\_simulation\_dimero\_trotter\_semplificato.ipynb} --
solo il richiesto dal relatore, ricerca dei parametri interattiva con
widget e persistenza.
\item \texttt{quantum\_simulation\_dimero\_trotter\_note.md} -- approfondimenti
tenuti fuori dal notebook completo (spiegazione ``non monocromatico'',
cronologia della selezione di $R_1$, convenzioni di segno).
\item \texttt{errata\_simmetrie\_dmi\_spin.tex/.pdf} e
\texttt{simmetrie\_dmi\_spin.tex/.pdf} (versione corretta integrale) --
correzione del fattore $2$/$\sqrt2$ negli elementi di matrice DM.
\item \texttt{dimero\_sintesi.tex/.pdf} (versione corretta) -- correzione
della norma dei commutatori in \S2.2.
\end{itemize}

\end{document}
